\section{Design Opportunities for Multisensory \& Multidimensional Biophilic Interactions}
\label{sec:ch4-design-opportunities}
We next synthesised Topics (TO1–TO7) and Themes (T1-T8) into five actionable design directions (DO1–DO5). We also explore how emerging technologies can reimagine biophilic experiences in smart environments, summarised in \autoref{tab:design_opportunities}. Although technology was not the primary focus of our interviews, experts highlighted its potential, particularly in themes related to art (T2), multi-species coexistence (T5), and embodied experiences (T7) (c.f. \autoref{sec:ch4-findings-a}). We used these discussions to support our design opportunities with technology enablers.


\paragraph{DO1: Designing for Emergent Natural Processes}
To foster deeper connection with nature, smart buildings must support evolving, self-regulating, and unpredictable qualities inherent in natural processes (see \autoref{sec:ch4-findings-thematic}). Unlike mechanical systems with repetitive patterns, natural processes evolve unpredictably, creating emergent properties---where small changes lead to transformative shifts~\cite{clement2004manifeste}. For example, ecosystems may adapt to support new species or behaviours, reflecting nature’s novelty and resilience~\cite{carscadden2023origins, radeloff2015rise}. Designing with this quality means embracing flux and variability---qualities rarely integrated into current smart building design (e.g.~\citet{economidou2021moving}'s vision of moving walls and talking floors). For IxD, this prompts the development of tools and methods that foreground emergence and unpredictability. For smart building researchers, it suggests investigating how such features influence occupants' comfort and emotional responses.

Realising emergent processes in smart buildings calls for technologies that enable unpredictability and self-regulation~\cite{clement2004manifeste}. This could involve repurposing the idea of ``third landscapes''---unmanaged, dedicated indoor areas where technology supports natural processes~\cite{clement2004manifeste}. \emph{Stochastic algorithms} in IoT systems might trigger environmental changes within probabilistic bounds, such as fluctuating light or airflow that reflects external wind variability~\cite{karapetyan2019multisensor}. \emph{Generative AI} could create evolving soundscapes in response to real-time environmental data—like wind speed, temperature, or occupancy, mirroring natural acoustic dynamism~\cite{parati2024using, ble2023generative}. \emph{Collaborative AI ecosystems}~\cite{koch2017design} might coordinate lighting, shading, or HVAC systems to behave in interdependent, temporally unfolding behaviours, generating a living quality within the building---similar to explorations in computational architecture~\cite{weber2024designing}.


\paragraph{DO2: Designing for Occupant Body Sensitivity}
This opportunity positions the human body as an interface and an active participant in biophilic experiences. By encouraging occupants to notice, understand, and reflect on their embodied sensations, buildings can foster deeper connections to the surrounding biophilic conditions  (c.f. Theme 7, \autoref{sec:ch4-findings-thematic}). For IxD, this opportunity challenges the prevailing paradigm of uniform sensory experiences in smart buildings, supporting first-person, body-centred interactions with biophilic conditions. 

Building on embodied urban HCI~\cite{galvao2024posthumanist}, technologies could amplify bodily awareness of ambient, subtle biophilic cues~\cite{tsaknaki2021feeling}. Haptic or thermal \emph{wearables} can translate airflow or temperature into tactile feedback. For instance, a wearable might vibrate or warm slightly to guide users to align themselves with airflow patterns in a room, encouraging a restorative state~\cite{sabinson2024every}. \emph{Smart building fa\c{c}ades could act as interactive sensory mediators}~\cite{srisuwan2022applications, rao2025smart}, responding to environmental changes such as ``sweating'' on hot days or during higher humidity levels before rain, connecting them with natural phenomena. \emph{Augmented Reality (AR)} can visualise real-time physiological data, such as heart rate or breathing, overlaying this information onto the user’s environment (similar to using AR for visualising indoor air pollution ~\cite{jarrahi2024enhancing}). An interface could, for instance, illustrate how proximity to plants or exposure to natural light impacts physiological responses, enabling reflection through a biophilic lens. Finally, \emph{digital journaling tools} like ``biophilic body maps'' could incorporate feminist and inclusive perspectives by allowing users to document and share how their identities and contexts shape their interactions with biophilic elements~\cite{turmo2023towards}. Aggregated body maps data from multiple users could reveal how groups engage with biophilic spaces, offering input for designing communal environments for maximising biophilic interactions.


\paragraph{DO3: Designing for Biophilic (Dis)Comfort}
Smart buildings should incorporate moments of \emph{(dis)comfort} to deepen human connections to nature (c.f. Theme 8, \autoref{sec:ch4-findings-thematic}). By integrating \emph{frictional} design elements inspired by biophilic surroundings~\cite{bixler1997nature}, IxD can create experiences that deliberately use discomfort to provoke reflection and inspire new understandings of the natural world~\cite{pezeu2013philosophy, rao2025smart}. 

IoT-enabled technology in flooring systems could simulate \emph{tactile interactions} by subtly cracking or shifting underfoot (e.g.~\citet{economidou2021moving}'s moving floors), symbolising ecosystem fragility and prompting occupants to reflect on climate impact. These might be augmented with ambient soundscapes, like the cracking of brittle surfaces to highlight ecosystem fragility~\cite{barclay2019acoustic}. We previously discussed \emph{biophilic (dis)comfort as a trade-off}, where occupants accept minor inconveniences for greater ecological or experiential benefits. Buildings could provide context for discomfort-inducing features, helping occupants interpret their experiences. For instance, AR overlays visualise indoor pollinator corridors with representations of bees or butterflies, paired with subtle sensory cues like vibrations or ambient buzzing sounds, encouraging reflection on shared habitats~\cite{mcclain2022using}. Discomfort could extend to \emph{experiences of risk and discernment}. As one expert noted, \textit{"Taste is a sense of caution... even in nature, we exercise a certain judgment."} Interactive installations could explore this cautious quality through scent or visual provocations that elicit hesitation and curiosity, prompting reflection on sensory boundaries (e.g. risk expressing design~\cite{disalvo2014making}). In workspaces, designers could introduce \emph{``productive biophilic discomfort''}~\cite{alpert2022productive}. For instance, smart lighting systems in light-scarce regions might dim briefly, mimicking cloud cover. While this momentary discomfort could disrupt focus, prompting occupants to seek natural light.


\paragraph{DO4: Designing for Biophilic Equity}
Through biophilic equity, smart buildings can address the systemic disparities in access to nature experiences and their benefits in indoor spaces, which has often been uneven (c.f. Themes 2 and 6, \autoref{sec:ch4-findings-thematic}). This necessitates prioritising public buildings and shared urban environments for biophilic design. For IxD and smart building researchers, biophilic equity offers opportunities to tackle the inequities through participatory technologies, therapeutic tools, and cultural storytelling. 

One potential application is \emph{digital platforms for participatory design}~\cite{arts2015digital}. Interactive apps or web platforms could enable community members to actively contribute to shaping biophilic spaces~\cite{arts2015digital, greene2020mobiles}. For instance, users could suggest plant species, vote on design features, or provide input that reflects their cultural and ecological values. \emph{Biophilic connection technologies} could also serve as therapeutic tools for marginalised groups such as displaced refugees or at-risk youth. Immersive virtual reality (VR) experiences could transport users to calming, nature-rich environments such as forests or ocean providing solace amidst hardships. \emph{Digital storytelling} through architectural features~\cite{burova2019promoting} in community centres or youth facilities could host interactive art installations or digital projections that celebrate the cultural and ecological histories of under-represented groups (like the co-design with displaced youth~\cite{fisher2016future} to develop wayfaring solutions). For example, public spaces could visualise the biodiversity of regions associated with displaced communities, fostering a shared ecological narrative and recognising their contributions to environmental stewardship.

\begingroup
\fontsize{7.4}{8.8}\selectfont
\setlength{\tabcolsep}{2.4pt}
\renewcommand{\arraystretch}{1.04}
\setlength{\LTleft}{0pt}
\setlength{\LTright}{0pt}

\begin{longtable}{@{}p{.22\linewidth} p{.33\linewidth} p{.11\linewidth} p{.24\linewidth}@{}}

\caption{A summary of design opportunities for supporting novel biophilic experiences in smart buildings. Each opportunity is summarised with a short description of how the smart building can be designed, corresponding themes, and three example technology possibilities. Theme 4 on multisensory and multidimensional experiences is not explicitly mentioned here as it is an underlying notion across all design opportunities. This is discussed further in \autoref{sec:ch4-discussions}.}
\label{tab:design_opportunities}\\

\toprule
\textbf{Design Opportunity} & \textbf{Description} & \textbf{Associated Theme(s)} & \textbf{Technology Examples (n=3)} \\
\midrule
\endfirsthead

\caption[]{A summary of design opportunities for supporting novel biophilic experiences in smart buildings, continued.}\\
\toprule
\textbf{Design Opportunity} & \textbf{Description} & \textbf{Associated Theme(s)} & \textbf{Technology Examples (n=3)} \\
\midrule
\endhead

\midrule
\multicolumn{4}{r}{\footnotesize Continued on next page}\\
\endfoot

\bottomrule
\endlastfoot

DO1 -- Designing for Emergent Natural Processes &
Support evolving, self-regulating, and unpredictable qualities of natural processes to create opportunities for ``new'' behaviours to emerge. &
T1, T2, T5 &
Stochastic IoT algorithms for environmental variability; generative AI for evolving soundscapes; collaborative AI ecosystems for generating complex building responses.\\
\midrule

DO2 -- Designing for Occupant Body Sensitivity &
Position the human body as an active participant in biophilic experiences by amplifying embodied awareness to surrounding biophilic conditions. &
T7 &
Haptic and thermal wearables for tactile feedback; AR to visualise physiological responses to natural elements; digital journaling tools for documenting interactions.\\
\midrule

DO3 -- Designing for Biophilic (Dis)Comfort &
Employ controlled discomfort in biophilic contexts to help occupants reflect on such moments and deepen human connections with nature. &
T8 &
AR overlays for pollinator corridors; soundscapes simulating ecosystem fragility; interactive installations exploring sensory boundaries around risk.\\
\midrule

DO4 -- Designing for Biophilic Equity &
Actively counter systemic disparities in access to nature-rich experiences by prioritising inclusive and participatory biophilic design that benefits underserved communities. &
T2, T6 &
Digital platforms for participatory design; VR experiences for immersive nature-rich environments; digital storytelling through projections or art installations.\\
\midrule

DO5 -- Designing for Relational Encounters with Nature &
Foster care, empathy, and shared responsibility by making nature's often invisible processes visible and encouraging collective interaction. &
T1, T3, T5 &
IoT sensors for ecological monitoring; AR for perspective-shifting experiences, e.g. plant's-eye view; AI-assisted wearables for ambient ecological awareness.\\

\end{longtable}
\endgroup

\paragraph{DO5 - Designing for Relational Encounters with Nature}
Smart buildings foster relational encounters by encouraging empathy, care, and mutual growth between humans and nature (c.f. Themes T3, T1, T5, \autoref{sec:ch4-findings-thematic}). It makes nature’s needs visible to promote awareness and responsibility~\cite{flanagan2018rewilding, li2022re}. Historical examples of Persian gardens embodied this by requiring active care from humans who maintained irrigation systems, creating ecosystems that rewarded engagement with fruit-bearing trees~\cite{khalilnezhad2016productive, khalilnezhad2024edible}. For researchers, designing for relational encounters with nature opens pathways for technologies fostering individual and collective empathy. It also calls for studying how group dynamics shape engagement within biophilic spaces.

Relational encounters with nature can be fostered by ``surfacing'' the presence and processes of ecological systems~\cite{flanagan2018rewilding, li2022re, steer2015growth}. This includes ``cues for care'' inviting reflection (discussed in T3 \autoref{sec:ch4-findings-thematic}).  Drawing from technology design in Relational-Cultural Theory (RCT)~\cite{comstock2008relational, kleespies2021connection, wiley2012empathy}, IoT sensors, for example, can track soil or light and translate data into interactive cues. For example, \emph{projection walls and ambient displays} could show environmental changes in shared spaces like lobbies or meeting rooms, supporting shared group responses to support biophilic elements~\cite{margariti2024evaluating}. \emph{Empathy-building technologies} could further enhance relational encounters by shifting perspectives from human to non-human~\cite{flanagan2018rewilding}. AR, for example, might simulate a ``plant’s-eye view'' of human actions~\cite{tracey2016does}, showing how behaviours like touching or crowding impact growth and health. \emph{Ecological interactions awareness} using AI-assisted wearables (e.g. similar to outdoor bird awareness bracelets~\cite{li2022re}) could be adapted for indoor environments, highlighting moving pollinators or subtle natural odours within biophilic spaces. 
